\documentclass[notes=show]{beamer}
\usepackage{mathpazo}
\usepackage{hyperref}
\usepackage{multimedia}
\usepackage{tcolorbox}
\usepackage{tikz}
\usetikzlibrary{shadows}
\usetheme{metropolis}
\usecolortheme{default}
\setbeamertemplate{itemize item}{%
    \begin{tikzpicture}
        \shade[ball color=black!100!yellow, preaction={fill=black,
opacity=.00,transform canvas={xshift=1mm,yshift=-1mm, yscale=0.5}}] (0,0) circle (0.6ex);
    \end{tikzpicture}
}
\setbeamertemplate{itemize subitem}{%
    \begin{tikzpicture}
        \shade[ball color=black!100!white, preaction={fill=black,
        opacity=.00,transform canvas={xshift=1mm,yshift=-1mm, yscale=0.5}}] (0,0) circle (0.6ex);
    \end{tikzpicture}
}
\setbeamercolor{frametitle}{bg=white}
\setbeamercolor{frametitle}{fg=black}
\setbeamercolor{background canvas}{bg=white}
\setbeamercolor{block body}{bg=mDarkTeal!30}
\setbeamercolor{block title}{bg=mDarkTeal,fg=black!2}
\setbeamertemplate{navigation symbols}{}
\setbeamertemplate{footline}[page number]
\newenvironment{stepenumerate}{\begin{enumerate}[<+->]}{\end{enumerate}}
\newenvironment{stepitemize}{\begin{itemize}[<+->]}{\end{itemize} }
\newenvironment{stepenumeratewithalert}{\begin{enumerate}[<+-| alert@+>]}{\end{enumerate}}
\newenvironment{stepitemizewithalert}{\begin{itemize}[<+-| alert@+>]}{\end{itemize} }

\graphicspath{{figures/}}

\begin{document}

\title{The Elusive Explanation for the Declining Labor Share}
\subtitle{}
\date{Gene Grossman \& Maurice Oberfield \bigskip \\
\textit{Annual Review of Economics}, 2022}
\author{}
\maketitle



%%%%%%%%%%%%%%%%%%%%%%%%%%%%%%%%%%%%%%%%%%%%%%%%%%%%%%%%%%%%%%%%%
\section{Introduction}
%%%%%%%%%%%%%%%%%%%%%%%%%%%%%%%%%%%%%%%%%%%%%%%%%%%%%%%%%%%%%%%%%

\begin{frame}{Introduction}
\begin{itemize}
\item Smith (1776) noted that the distribution of national income to wages, rents, and profits was closely related to inequality. \medskip
\item Kaldor (1961) named the constancy of the labor share as the first of his stylized facts about economic growth. \medskip
\item After more than a century of relative stability, the labor share has declined recently. \medskip
\item This has lead to revival of interest into the labor share.
\end{itemize}
\end{frame}

\AtBeginSection[ ]
{
\begin{frame}{Outline}
    \tableofcontents[currentsection]
\end{frame}
}

%%%%%%%%%%%%%%%%%%%%%%%%%%%%%%%%%%%%%%%%%%%%%%%%%%%%%%%%%%%%%%%%%
\section{Measuring the decline in the labor share}
%%%%%%%%%%%%%%%%%%%%%%%%%%%%%%%%%%%%%%%%%%%%%%%%%%%%%%%%%%%%%%%%%

\begin{frame}{The labor share in the US}
\begin{figure}[p!]
 \includegraphics[width=.9\textwidth]{figure1.pdf} 
\end{figure}
\end{frame}

\begin{frame}{Various measures of labor shares}
\begin{itemize}
\item Various measures of labor shares depending on: \medskip
\begin{enumerate}
\item Gross versus net value added \medskip
\item Treatment of Intellectual Property (IP) \medskip
\item Treatment of self-employment and entrepreneurial income \medskip
\item Treatment of owner-occupied housing \medskip
\item Treatment of factorless income
\end{enumerate} \medskip
\item Nonetheless, there is a consensus that the labor share today is below its level of 1970s and 1980s in the US. \medskip
\item Same is true for many other countries.
\end{itemize}
\end{frame}

\begin{frame}{The global labor share (Karabarbounis \& Neiman 2014)}
\begin{figure}[p!]
 \includegraphics[width=.7\textwidth, angle=-90]{figures/figure-KN-1.pdf} 
\end{figure}
\end{frame}

%%%%%%%%%%%%%%%%%%%%%%%%%%%%%%%%%%%%%%%%%%%%%%%%%%%%%%%%%%%%%%%%%
\section{Some pitfalls using cross-sectional data}
%%%%%%%%%%%%%%%%%%%%%%%%%%%%%%%%%%%%%%%%%%%%%%%%%%%%%%%%%%%%%%%%%

\begin{frame}{Defining the labor share}
\begin{itemize}
\item Assume a CRS production function $Y = F(K,L; \tau)$. \medskip
\item Parameter $\tau$ influences relative factor demands. \medskip
\item The unit cost function is $c(r,w; \tau)$. \medskip
\item The labor share is given by:
\begin{equation*}
    \theta_{L} \equiv \frac{wL}{Y} = \frac{w c_{w}(r,w;\tau)Y}{Y} = w c_{w}(r,w;\tau)    
\end{equation*}
\item The capital share is given by $\theta_{K} = r c_{r}(r,w;\tau)$
\end{itemize}
\end{frame}

\begin{frame}{Changes in the labor share}
Differentiating $ \ln(\theta_{L} / \theta_{K}) $ with respect to $\tau$ analyzes demand-driven changes in the labor share:
\begin{align*}
    \frac{d \ln(\theta_{L} / \theta_{K})}{d \tau} & = \frac{d \ln(wc_{w} / rc_{r})}{d \tau} \nonumber \\
    & = \frac{\partial \ln(c_{w} / c_{r})}{\partial \tau} + \left[ \frac{\partial \ln(wc_{w} / rc_{r})}{\partial \ln(r / w )} \right] \frac{d \ln(r / w)}{d \tau} \nonumber \\
     & = \frac{\partial \ln(c_{w} / c_{r})}{\partial \tau} + \left[ \frac{\partial \ln(w / r)}{\partial \ln(r / w ) } + \frac{\partial \ln(L / K)}{\partial \ln(r / w ) } \right] \frac{d \ln(r / w)}{d \tau} \nonumber \\
    & = \frac{\partial \ln(L / K) }{\partial \tau} + [ \sigma - 1 ] \frac{d \ln(r / w) }{d \tau} \tag{1}
\end{align*}
with $\sigma > 0 $ the elasticity of substitution between $K$ and $L$.
\end{frame}

\begin{frame}{Changes in the labor share}
\begin{itemize}
\item Decomposition of demand-driven changes in the labor share:
\begin{equation*}
    \frac{d \ln(\theta_{L} / \theta_{K})}{d \tau} = \frac{\partial \ln(L / K) }{\partial \tau} + [ \sigma - 1 ] \frac{d \ln(r / w) }{d \tau} \tag{1}
\end{equation*}
\item The labor share changes because of a \textit{shift in relative factor demands} (i.e. a change in relative factor quantities for given relative factor prices) and a \textit{change in relative factor prices}. \medskip
\item $ \sigma  \geq 1$ plays a key role. \medskip
\item Equation (1) is general and must hold for all $ d \tau $.
\end{itemize}
\end{frame}

\begin{frame}{Example: BGP and Harrod-neutral technological progress}
\begin{itemize}
\item In growth models, a BGP is defined by constant growth in $Y/L$ and constant $K/Y$, $r$ and $ \theta_{L}/\theta_{K}$ (Kaldor 1963). \medskip  
\item In Solow growth models, Kaldor facts require tech progress is Harrod-neutral (purely labor-augmenting) (Uzawa 1961). \medskip
\item Wage grows at constant rate $ \tau $: $ d \ln(w) = \tau $. \medskip
\item Also: $ c_{w}(r,w; \tau) = c_{w} (1, \omega; \tau) = e^{-\tau} \tilde{c}_{w} (1, \omega e^{-\tau}) $ and  $ c_{r}(r,w; \tau) = c_{r} (1, \omega; \tau) = \tilde{c}_{r} (1, \omega e^{-\tau}) $. \medskip
\item Given this, is equation (1) consistent with $ d \ln( \theta_{L} / \theta_{K})=0 $?
\end{itemize}
\end{frame}

\begin{frame}{Example: BGP and Harrod-neutral technological progress}
\begin{itemize}
\item When deriving equation (1), we used that:
\begin{align*}
    \sigma - 1 & = \frac{\partial \ln(rc_{r}(r,w;\tau) / wc_{w}(r,w;\tau))}{\partial \ln(w / r )} \\
    & = \frac{\partial \ln(c_{r}(r,w;\tau) / \omega c_{w}(r,w;\tau))}{\partial \ln(\omega)} \\
    & = \frac{\partial \ln(c_{r}(1,\omega;\tau) / \omega c_{w}(1, \omega;\tau))}{\partial \ln(\omega)}
\end{align*}
\item Kaldor facts and Harrod-neutral technological progress imply:
\begin{align*}
    \sigma - 1 & = \frac{\partial \ln(\tilde{c}_{r}(1,\omega e^{-\tau}) / \omega e^{-\tau} \tilde{c}_{w}(1, \omega e^{-\tau}))}{\partial \ln(\omega)}  \\
    & = \frac{\partial \ln(\tilde{c}_{w}(1,\omega e^{-\tau}) / \tilde{c}_{r}(1, \omega e^{-\tau}))}{\partial \tau} 
\end{align*}
\end{itemize}
\end{frame}

\begin{frame}{Pitfalls when estimating changes in factor shares}
\begin{itemize}
\item Let $d \tau_{i}$ be a shock in firm or region $i$ and estimate: \medskip
\begin{equation*}
    d \ln(\theta_{L,i} / (1-\theta_{L,i}) = \alpha + \beta d \tau_{i} + \epsilon_{i} \tag{2}
\end{equation*}
with $d \tau_{i} $ a shock measure (e.g. tech progress). \medskip
\item If relative factor prices don't vary across $i$, $\hat{\beta} $ captures the partial equilibrium impact of $ d \tau_{i}$ on shifts in relative factor demands (e.g. $\beta= 0 $ $\Rightarrow $ tech progress is Hicks-neutral).
\medskip
\item If relative factor demand-shifts and prices vary across $i$, $\hat{\beta} $ captures the general equilibrium impact of $ d \tau_{i}$ (e.g. $\beta= 0 $ $\Rightarrow $ tech progress is Harrod-neutral).
\end{itemize}
\end{frame}

%%%%%%%%%%%%%%%%%%%%%%%%%%%%%%%%%%%%%%%%%%%%%%%%%%%%%%%%%%%%%%%%%
\section{Technical change}
%%%%%%%%%%%%%%%%%%%%%%%%%%%%%%%%%%%%%%%%%%%%%%%%%%%%%%%%%%%%%%%%%

%%%%%%%%%%%%%%%%%%%%%%%%%%%%%%%%%%%%%%%%%%%%%%%%%%%%%%%%%%%%%%%%%
\subsection*{Investment-specific technological change}
%%%%%%%%%%%%%%%%%%%%%%%%%%%%%%%%%%%%%%%%%%%%%%%%%%%%%%%%%%%%%%%%%

\begin{frame}
	\centering
	\textbf{Investment-specific technological change}
\end{frame}

\begin{frame}{Decline in price of investment goods}
\begin{figure}[p!]
 \includegraphics[width=.8\textwidth]{figure-KN-7.pdf} 
\end{figure}
\end{frame}

\begin{frame}{Investment-specific technological change}
\begin{itemize}
\item Assume tech progress decreases the relative price of investment goods $R_{i}$. \medskip
\item Karabarbounis \& Neiman (2014) use data for 59 countries for 1975-2012 and a general equilibrium equation:
\begin{equation*}
    \frac{\theta_{L,i}}{1 - \theta_{L,i}} d \ln(\theta_{L,i}) = \gamma + [\sigma - 1] d \ln(R_{i}) + \epsilon_{i} \tag{KN-19}
\end{equation*}
with an estimate $\hat{\sigma} = 1.25$. \medskip
\item A decrease in $R_{i}$ induces a substitution from labor to capital to such extent that the labor share falls.
\end{itemize}
\end{frame}

\begin{frame}{Labor share and price of investment goods}
\begin{figure}[p!]
 \includegraphics[width=.7\textwidth]{figure-KN-8.pdf} 
\end{figure}
\end{frame}

%%%%%%%%%%%%%%%%%%%%%%%%%%%%%%%%%%%%%%%%%%%%%%%%%%%%%%%%%%%%%%%%%
\subsection*{Robot adoption}
%%%%%%%%%%%%%%%%%%%%%%%%%%%%%%%%%%%%%%%%%%%%%%%%%%%%%%%%%%%%%%%%%

\begin{frame}
	\centering
	\textbf{Robot adoption}
\end{frame}

\begin{frame}{Automation and creation of labor tasks}
\begin{itemize}
\item Automation (\textit{creation}) of labor tasks assumes that capital (\textit{labor}) directly substitutes for some labor (\textit{capital}) tasks. \medskip
\item Acemoglu \& Restrepo (2019) use data for US sectors for 1947-1987 and estimate:
\begin{equation*}
    \frac{1}{1 - \theta_{L,i}} d \ln(\theta_{L,i}) = \frac{d \ln(\Gamma_{i})}{1- \Gamma_{i}} + [\sigma -1] d \ln(\frac{r_{i}}{w_{i}}) \tag{AR-A10}
\end{equation*}
with $ d \ln(\Gamma_{i}) / (1- \Gamma_{i})$ shifts in relative labor demand due to automation (-) or creation (+) of labor tasks. \medskip
\item Assuming $\sigma = 0.8$, they impute $ d \ln(\Gamma_{i}) / (1- \Gamma_{i}) $ and find that automation has decreased $ \theta_{L} $ after 1987.
\end{itemize}
\end{frame}

\begin{frame}{Shifts in relative labor demand (Figure 5)}
\begin{figure}[p!]
 \includegraphics[width=\textwidth]{figure-AR-5B.pdf} 
\end{figure}
\end{frame}

\begin{frame}{Shifts in relative labor demand and robot adoption (Figure 4A)}
\begin{figure}[p!]
 \includegraphics[width=\textwidth]{figure-AR-A4.pdf} 
\end{figure}
\end{frame}

%%%%%%%%%%%%%%%%%%%%%%%%%%%%%%%%%%%%%%%%%%%%%%%%%%%%%%%%%%%%%%%%%
\subsection*{The rise of superstar firms}
%%%%%%%%%%%%%%%%%%%%%%%%%%%%%%%%%%%%%%%%%%%%%%%%%%%%%%%%%%%%%%%%%

\begin{frame}
	\centering
	\textbf{The rise of superstar firms}
\end{frame}

\begin{frame}{The rise of superstar firms (Autor et al. 2020)}
\begin{itemize}
\item The sales-weighted labor share decreases if superstar firms become more important and have lower labor shares. \medskip
\item Autor et al. (2020) decompose the sales-weighted labor share:
\begin{equation*}
    S = \sum \omega_{i} S_{i} = \bar{S} + \sum (\omega_{i} - \bar{\omega})(S_{i}-\bar{S}) \tag{ADKPR-3}
\end{equation*}
with $ \omega_{i} \equiv P{i}Y_{i}/PY$, $\bar{S}$ the unweighted mean labor share, and $\bar{\omega}$ the unweighted mean value-added share. \medskip
\item Change in the labor share can be written as:
\begin{equation*}
    \Delta S = \Delta \bar{S} + \Delta \left[ \sum (\omega_{i} - \bar{\omega})(S_{i}-\bar{S}) \right] \tag{ADKPR-4}
\end{equation*}
\end{itemize}
\end{frame}

\begin{frame}{The rise of superstar firms decreases the labor share}
\begin{figure}[p!]
 \includegraphics[width=0.7\textwidth]{figure-ADKPR-7.pdf} 
\end{figure}
\end{frame}

%%%%%%%%%%%%%%%%%%%%%%%%%%%%%%%%%%%%%%%%%%%%%%%%%%%%%%%%%%%%%%%%%
\section{Globalization and the rise of China}
%%%%%%%%%%%%%%%%%%%%%%%%%%%%%%%%%%%%%%%%%%%%%%%%%%%%%%%%%%%%%%%%%

\begin{frame}{Globalization and the rise of China}
\begin{itemize}
\item When do imported inputs decrease the domestic labor share? \medskip 
\item Elsby, Hobijn \& Sahin (2013) use data for US sectors for 1993-2010 and the equation:
\begin{align*}
    \frac{1}{1-\theta_{L,i}} \frac{d \ln(\theta_{L,i})}{d \ln(M_{i})} & = - s_{M,i} \left( \rho_{LM} - \rho_{KM} \right) \tag{EHS-11} \\
    & = - s_{M,i} \left( \frac{\sigma_{LM}-1}{\sigma_{LM}} - \frac{\sigma_{KM}-1}{\sigma_{KM}} \right)  
\end{align*}
with $s_{M,j} \equiv M_{i} / (P_{i}Y_{i})$. \medskip
\item If $ \sigma_{LM} > \sigma_{KM} $, imported inputs will induce greater substitution away from labor than from capital and $ \theta_{L,i}$ decreases.
\end{itemize}
\end{frame}

\begin{frame}{The labor share and import competition across industries}
\begin{figure}[p!]
 \includegraphics[width=.8\textwidth]{figure-EHS-12.pdf} 
\end{figure}
\end{frame}

%%%%%%%%%%%%%%%%%%%%%%%%%%%%%%%%%%%%%%%%%%%%%%%%%%%%%%%%%%%%%%%%%
\section{Increased product market power}
%%%%%%%%%%%%%%%%%%%%%%%%%%%%%%%%%%%%%%%%%%%%%%%%%%%%%%%%%%%%%%%%%

\begin{frame}{The rise of markups (De Loecker, Eeckhout \& Unger 2020)}
\begin{figure}[p!]
 \includegraphics[width=\textwidth]{figure-LEU-1.pdf} 
\end{figure}
\end{frame}

\begin{frame}{Monopoly and price markups}
\begin{itemize}
\item Why do higher markups lower the labor share? \medskip
\item Assuming $ Y= F(K,L)$ and inverse product demand given by $p=P(Y)$, profit maximization implies:
\begin{equation*}
    \rho F_{L} = w
\end{equation*}
with $ \rho \equiv p+YP_{Y}$ marginal revenue.
\item Rearranging terms and multiplying through by $L/Y$:
\begin{equation*}
    \epsilon_{YL} \equiv \frac{LF_{L}}{Y} = \frac{wL}{\rho Y} = \frac{p}{\rho} \frac{wL}{pY} = \frac{p}{\rho} \theta_{L}
\end{equation*}
with $  \epsilon_{YL}$ the elasticity of $Y$ with respect to $L$.
\end{itemize}
\end{frame}

\begin{frame}{Markups decrease factor shares}
\begin{itemize}
\item Equating marginal revenue to marginal cost, $\rho = MC$:
\begin{equation*}
    \epsilon_{YL} = \frac{p}{\rho} \theta_{L}  = \frac{p}{MC} \theta_{L} = \mu \theta_{L} \tag{4}
\end{equation*}
with $  \mu $ the markup of output price over marginal cost. \medskip
\item The analogous expression for capital is:
\begin{equation*}
    \epsilon_{YK} = \frac{p}{\rho} \theta_{K} = \frac{p}{MC} \theta_{K} = \mu \theta_{K}
\end{equation*}
\item For given $\epsilon_{YL} $ and $\epsilon_{YK} $, an increase in $\mu$ decreases $\theta_{L} $ and $\theta_{K} $.
\end{itemize}
\end{frame}

\begin{frame}{Markups increase the profit share}
\begin{itemize}
\item Summing $ \theta_{L} = \epsilon_{YL} / \mu $ and $ \theta_{K} = \epsilon_{YK} / \mu  $ gives:
\begin{equation*}
    \theta_{L} + \theta_{K}  = \frac{\epsilon_{YL} + \epsilon_{YK}}{\mu} = 1 - \theta_{\pi} \tag{5}
\end{equation*}
with $\theta_{\pi} $ the profit share and $\epsilon_{YL} + \epsilon_{YK} $ return to scale. \medskip
\item For given $ \epsilon_{YL} + \epsilon_{YK}$, an increase in $  \mu $ increases $ \theta_{\pi}$. \medskip
\item Alternatively, we can write for $j=L,K$:
\begin{equation*}
    \theta_{j} = \frac{\epsilon_{Yj}}{\epsilon_{YL} + \epsilon_{YK}}[1- \theta_{\pi}]  \tag{6}
\end{equation*}
showing that an increase in $\theta_{\pi} $ decreases $\theta_{L} $ and $\theta_{K} $.
\end{itemize}
\end{frame}

\begin{frame}{Estimating and interpreting rising markups}
\begin{itemize}
\item De Loecker, Eeckhout \& Unger (2020) estimate firm-level markups using a composite of variable inputs:
\begin{equation*}
    \hat{\mu}_{it} = \frac{\hat{\epsilon}_{YV,t}}{\hat{\theta}_{V,it}} 
\end{equation*}
with $ \epsilon_{YV,t}$ estimated using Olley-Pakes technique. \medskip
\item The sales-weighted average markup in year $t$ is given by:
\begin{equation*}
    \mathbb{E}_{t}{[\hat{\mu}_{i,t}]} = \sum_{i} \frac{p_{it}y_{it}}{P_{t}Y_{t}} \frac{\hat{\epsilon}_{YV,t}}{\hat{\theta}_{V,it}} 
\end{equation*}
\item Measures have been critized and rising markups is a proximate (not primitive) cause of fall in $ \theta_{j}$.
\end{itemize}
\end{frame}

%%%%%%%%%%%%%%%%%%%%%%%%%%%%%%%%%%%%%%%%%%%%%%%%%%%%%%%%%%%%%%%%%
\section{Declining market power of workers}
%%%%%%%%%%%%%%%%%%%%%%%%%%%%%%%%%%%%%%%%%%%%%%%%%%%%%%%%%%%%%%%%%

%%%%%%%%%%%%%%%%%%%%%%%%%%%%%%%%%%%%%%%%%%%%%%%%%%%%%%%%%%%%%%%%%
\subsection*{The decline in union membership}
%%%%%%%%%%%%%%%%%%%%%%%%%%%%%%%%%%%%%%%%%%%%%%%%%%%%%%%%%%%%%%%%%

\begin{frame}
	\centering
	\textbf{The decline in union membership}
\end{frame}

\begin{frame}{Declining union membership}
\begin{itemize}
\item When does de-unionisation decrease the labor share? \medskip
\item Bentolila \& Saint-Paul (2003) use CES production to get:
\begin{equation*}
    \theta_{L,i} = 1 - \alpha k_{i}^{\frac{\sigma - 1}{\sigma}} \tag{BS-10}
\end{equation*}
with $ k_{i} \equiv K_{i} / Y_{i} $ and $ \alpha$ capital's distribution parameter. \medskip
\item A decrease in the union-wage premium makes firms substitute labor for capital which decreases $k_{i}$. \medskip
\item If $\sigma < 1 $, de-unionisation decreases $\theta_{L,i}$.
\end{itemize}
\end{frame}

\begin{frame}{Labor share and labor rents (Stansbury \& Summers 2020)}
\begin{figure}[p!]
 \includegraphics[width=.9\textwidth]{figure-SS-9.pdf} 
\end{figure}
\end{frame}

%%%%%%%%%%%%%%%%%%%%%%%%%%%%%%%%%%%%%%%%%%%%%%%%%%%%%%%%%%%%%%%%%
\subsection*{Increased concentration of labor demand}
%%%%%%%%%%%%%%%%%%%%%%%%%%%%%%%%%%%%%%%%%%%%%%%%%%%%%%%%%%%%%%%%%

\begin{frame}
	\centering
	\textbf{Increased concentration of labor demand}
\end{frame}

\begin{frame}{Monopsony and wage markdowns}
\begin{itemize}
\item Why do higher wage markdowns lower the labor share? \medskip
\item Assuming $ Y= F(K,L)$ and inverse labor supply $w=W(L)$, profit maximization implies:
\begin{equation*}
    p F_{L} = \tilde{\mu} w
\end{equation*}
with $ \tilde{\mu} \equiv 1+ \epsilon_{Lw}^{-1} $ and $ \epsilon_{Lw}^{-1}$ the elasticity of $W(L)$.
\item Rearranging terms and multiplying through by $L/Y$:
\begin{equation*}
    \epsilon_{YL} \equiv \frac{LF_{L}}{Y} = \tilde{\mu} \frac{wL}{p Y}  = \tilde{\mu} \theta_{L}
\end{equation*}
such that an increase in $\tilde{\mu}$ decreases $\theta_{L}$ for given $\epsilon_{YL}$.
\end{itemize}
\end{frame}

\begin{frame}{ Wages and employer concentration (Azar et al. 2020)}
\begin{figure}[p!]
 \includegraphics[width=.9\textwidth]{figure-AMS-5.pdf} 
\end{figure}
\end{frame}

%%%%%%%%%%%%%%%%%%%%%%%%%%%%%%%%%%%%%%%%%%%%%%%%%%%%%%%%%%%%%%%%%
\section{Conclusions}
%%%%%%%%%%%%%%%%%%%%%%%%%%%%%%%%%%%%%%%%%%%%%%%%%%%%%%%%%%%%%%%%%

\begin{frame}{Conclusions}
\begin{itemize}
\item After remaining constant for a long time, the labor share has been declining in the last decades. \medskip
\item Several reasons have been documented in the literature: \medskip
\begin{itemize}
    \item Technological progress and the rise of superstar firms. \medskip
    \item Increased imported inputs. \medskip
    \item De-unionisation. \medskip
    \item Increased price markups and wage markdowns. \medskip
\end{itemize}
\item There is overcounting because explanations are not exclusive and not all primitive drivers of change. 
\end{itemize}
\end{frame}

\begin{frame}{Further readings}
\footnotesize
\begin{enumerate}
\item Karabarbounis, L., and B. Neiman (2014b). “The Global Decline of the Labor Share”. \textit{Quarterly Journal of Economics}. 129(1): 61–103.
\item Elsby, M., B. Hobijn, and A. Sahin (2013). “The Decline of the U.S. Labor Share”. \textit{Brookings Papers on Economic Activity}. 44: 1–52.
\item Acemoglu, D., and P. Restrepo (2019). “Automation and New Tasks: How Technology Displaces and Reinstates Labor”. \textit{Journal of Economic Perspectives}. 33:2: 3-30.
\item Stansbury, A., and L. Summers (2020). “The Declining Worker Hypothesis: An Explanation for the Recent Evolution of the American Economy”. \textit{Brookings Papers on Economic Activity}. Spring 2020
\item Autor, D., Dorn, D., Katz, L., Patterson, C., and J. Van Reenen (2020). “The Fall of the Labor Share and the Rise of Superstar Firms”. \textit{ Quarterly Journal of Economics}. 135(2): 645-709.
\item L. Karabarbounis (2023). “Perspectives on the Labor Share”. \textit{NBER Working Paper} 31854. November 2023.
\end{enumerate}
\end{frame}

\end{document}
